\IfFileExists{stacks-project.cls}{%
\documentclass{stacks-project}
}{%
\documentclass{amsart}
}

% For dealing with references we use the comment environment
\usepackage{verbatim}
\newenvironment{reference}{\comment}{\endcomment}
%\newenvironment{reference}{}{}
\newenvironment{slogan}{\comment}{\endcomment}
\newenvironment{history}{\comment}{\endcomment}

% For commutative diagrams we use Xy-pic
\usepackage[all]{xy}

% We use 2cell for 2-commutative diagrams.
\xyoption{2cell}
\UseAllTwocells

% We use multicol for the list of chapters between chapters
\usepackage{multicol}

% This is generally recommended for better output
\usepackage{lmodern}
\usepackage[T1]{fontenc}

% For cross-file-references
\usepackage{xr-hyper}

% Package for hypertext links:
\usepackage{hyperref}

% For any local file, say "hello.tex" you want to link to please
% use \externaldocument[hello-]{hello}
\externaldocument[introduction-]{introduction}
\externaldocument[conventions-]{conventions}
\externaldocument[sets-]{sets}
\externaldocument[categories-]{categories}
\externaldocument[topology-]{topology}
\externaldocument[sheaves-]{sheaves}
\externaldocument[sites-]{sites}
\externaldocument[stacks-]{stacks}
\externaldocument[fields-]{fields}
\externaldocument[algebra-]{algebra}
\externaldocument[brauer-]{brauer}
\externaldocument[homology-]{homology}
\externaldocument[derived-]{derived}
\externaldocument[simplicial-]{simplicial}
\externaldocument[more-algebra-]{more-algebra}
\externaldocument[smoothing-]{smoothing}
\externaldocument[modules-]{modules}
\externaldocument[sites-modules-]{sites-modules}
\externaldocument[injectives-]{injectives}
\externaldocument[cohomology-]{cohomology}
\externaldocument[sites-cohomology-]{sites-cohomology}
\externaldocument[dga-]{dga}
\externaldocument[dpa-]{dpa}
\externaldocument[sdga-]{sdga}
\externaldocument[hypercovering-]{hypercovering}
\externaldocument[schemes-]{schemes}
\externaldocument[constructions-]{constructions}
\externaldocument[properties-]{properties}
\externaldocument[morphisms-]{morphisms}
\externaldocument[coherent-]{coherent}
\externaldocument[divisors-]{divisors}
\externaldocument[limits-]{limits}
\externaldocument[varieties-]{varieties}
\externaldocument[topologies-]{topologies}
\externaldocument[descent-]{descent}
\externaldocument[perfect-]{perfect}
\externaldocument[more-morphisms-]{more-morphisms}
\externaldocument[flat-]{flat}
\externaldocument[groupoids-]{groupoids}
\externaldocument[more-groupoids-]{more-groupoids}
\externaldocument[etale-]{etale}
\externaldocument[chow-]{chow}
\externaldocument[intersection-]{intersection}
\externaldocument[pic-]{pic}
\externaldocument[weil-]{weil}
\externaldocument[adequate-]{adequate}
\externaldocument[dualizing-]{dualizing}
\externaldocument[duality-]{duality}
\externaldocument[discriminant-]{discriminant}
\externaldocument[derham-]{derham}
\externaldocument[local-cohomology-]{local-cohomology}
\externaldocument[algebraization-]{algebraization}
\externaldocument[curves-]{curves}
\externaldocument[resolve-]{resolve}
\externaldocument[models-]{models}
\externaldocument[functors-]{functors}
\externaldocument[equiv-]{equiv}
\externaldocument[pione-]{pione}
\externaldocument[etale-cohomology-]{etale-cohomology}
\externaldocument[proetale-]{proetale}
\externaldocument[relative-cycles-]{relative-cycles}
\externaldocument[more-etale-]{more-etale}
\externaldocument[trace-]{trace}
\externaldocument[crystalline-]{crystalline}
\externaldocument[spaces-]{spaces}
\externaldocument[spaces-properties-]{spaces-properties}
\externaldocument[spaces-morphisms-]{spaces-morphisms}
\externaldocument[decent-spaces-]{decent-spaces}
\externaldocument[spaces-cohomology-]{spaces-cohomology}
\externaldocument[spaces-limits-]{spaces-limits}
\externaldocument[spaces-divisors-]{spaces-divisors}
\externaldocument[spaces-over-fields-]{spaces-over-fields}
\externaldocument[spaces-topologies-]{spaces-topologies}
\externaldocument[spaces-descent-]{spaces-descent}
\externaldocument[spaces-perfect-]{spaces-perfect}
\externaldocument[spaces-more-morphisms-]{spaces-more-morphisms}
\externaldocument[spaces-flat-]{spaces-flat}
\externaldocument[spaces-groupoids-]{spaces-groupoids}
\externaldocument[spaces-more-groupoids-]{spaces-more-groupoids}
\externaldocument[bootstrap-]{bootstrap}
\externaldocument[spaces-pushouts-]{spaces-pushouts}
\externaldocument[spaces-chow-]{spaces-chow}
\externaldocument[groupoids-quotients-]{groupoids-quotients}
\externaldocument[spaces-more-cohomology-]{spaces-more-cohomology}
\externaldocument[spaces-simplicial-]{spaces-simplicial}
\externaldocument[spaces-duality-]{spaces-duality}
\externaldocument[formal-spaces-]{formal-spaces}
\externaldocument[restricted-]{restricted}
\externaldocument[spaces-resolve-]{spaces-resolve}
\externaldocument[formal-defos-]{formal-defos}
\externaldocument[defos-]{defos}
\externaldocument[cotangent-]{cotangent}
\externaldocument[examples-defos-]{examples-defos}
\externaldocument[algebraic-]{algebraic}
\externaldocument[examples-stacks-]{examples-stacks}
\externaldocument[stacks-sheaves-]{stacks-sheaves}
\externaldocument[criteria-]{criteria}
\externaldocument[artin-]{artin}
\externaldocument[quot-]{quot}
\externaldocument[stacks-properties-]{stacks-properties}
\externaldocument[stacks-morphisms-]{stacks-morphisms}
\externaldocument[stacks-limits-]{stacks-limits}
\externaldocument[stacks-cohomology-]{stacks-cohomology}
\externaldocument[stacks-perfect-]{stacks-perfect}
\externaldocument[stacks-introduction-]{stacks-introduction}
\externaldocument[stacks-more-morphisms-]{stacks-more-morphisms}
\externaldocument[stacks-geometry-]{stacks-geometry}
\externaldocument[moduli-]{moduli}
\externaldocument[moduli-curves-]{moduli-curves}
\externaldocument[examples-]{examples}
\externaldocument[exercises-]{exercises}
\externaldocument[guide-]{guide}
\externaldocument[desirables-]{desirables}
\externaldocument[coding-]{coding}
\externaldocument[obsolete-]{obsolete}
\externaldocument[fdl-]{fdl}
\externaldocument[index-]{index}

% Theorem environments.
%
\theoremstyle{plain}
\newtheorem{theorem}[subsection]{Theorem}
\newtheorem{proposition}[subsection]{Proposition}
\newtheorem{lemma}[subsection]{Lemma}

\theoremstyle{definition}
\newtheorem{definition}[subsection]{Definition}
\newtheorem{example}[subsection]{Example}
\newtheorem{exercise}[subsection]{Exercise}
\newtheorem{situation}[subsection]{Situation}

\theoremstyle{remark}
\newtheorem{remark}[subsection]{Remark}
\newtheorem{remarks}[subsection]{Remarks}

\numberwithin{equation}{subsection}

% Macros
%
\def\lim{\mathop{\mathrm{lim}}\nolimits}
\def\colim{\mathop{\mathrm{colim}}\nolimits}
\def\Spec{\mathop{\mathrm{Spec}}}
\def\Hom{\mathop{\mathrm{Hom}}\nolimits}
\def\Ext{\mathop{\mathrm{Ext}}\nolimits}
\def\SheafHom{\mathop{\mathcal{H}\!\mathit{om}}\nolimits}
\def\SheafExt{\mathop{\mathcal{E}\!\mathit{xt}}\nolimits}
\def\Sch{\mathit{Sch}}
\def\Mor{\mathop{\mathrm{Mor}}\nolimits}
\def\Ob{\mathop{\mathrm{Ob}}\nolimits}
\def\Sh{\mathop{\mathit{Sh}}\nolimits}
\def\NL{\mathop{N\!L}\nolimits}
\def\CH{\mathop{\mathrm{CH}}\nolimits}
\def\proetale{{pro\text{-}\acute{e}tale}}
\def\etale{{\acute{e}tale}}
\def\QCoh{\mathit{QCoh}}
\def\Ker{\mathop{\mathrm{Ker}}}
\def\Im{\mathop{\mathrm{Im}}}
\def\Coker{\mathop{\mathrm{Coker}}}
\def\Coim{\mathop{\mathrm{Coim}}}

% Boxtimes
%
\DeclareMathSymbol{\boxtimes}{\mathbin}{AMSa}{"02}

%
% Macros for moduli stacks/spaces
%
\def\QCohstack{\mathcal{QC}\!\mathit{oh}}
\def\Cohstack{\mathcal{C}\!\mathit{oh}}
\def\Spacesstack{\mathcal{S}\!\mathit{paces}}
\def\Quotfunctor{\mathrm{Quot}}
\def\Hilbfunctor{\mathrm{Hilb}}
\def\Curvesstack{\mathcal{C}\!\mathit{urves}}
\def\Polarizedstack{\mathcal{P}\!\mathit{olarized}}
\def\Complexesstack{\mathcal{C}\!\mathit{omplexes}}
% \Pic is the operator that assigns to X its picard group, usage \Pic(X)
% \Picardstack_{X/B} denotes the Picard stack of X over B
% \Picardfunctor_{X/B} denotes the Picard functor of X over B
\def\Pic{\mathop{\mathrm{Pic}}\nolimits}
\def\Picardstack{\mathcal{P}\!\mathit{ic}}
\def\Picardfunctor{\mathrm{Pic}}
\def\Deformationcategory{\mathcal{D}\!\mathit{ef}}
\usepackage{cleveref}
\crefname{originalsection}{original section}{original sections}
\Crefname{originalsection}{Original section}{Original sections}
\makeatletter
\newlabel{algebra-section-going-up}{{00HU}{source}{}{sourcepointer-1}{}}
\newlabel{algebra-section-going-up@cref}{{[originalsection][0][]00HU}{[1][0][]source}{}{}{}}
\makeatother
\crefname{originallemma}{original lemma}{original lemmas}
\Crefname{originallemma}{Original lemma}{Original lemmas}
\makeatletter
\newlabel{morphisms-lemma-etale-flat}{{02GS}{source}{}{sourcepointer-2}{}}
\newlabel{morphisms-lemma-etale-flat@cref}{{[originallemma][0][]02GS}{[1][0][]source}{}{}{}}
\makeatother
\crefname{originallemma}{original lemma}{original lemmas}
\Crefname{originallemma}{Original lemma}{Original lemmas}
\makeatletter
\newlabel{morphisms-lemma-etale-over-field}{{02GL}{source}{}{sourcepointer-3}{}}
\newlabel{morphisms-lemma-etale-over-field@cref}{{[originallemma][0][]02GL}{[1][0][]source}{}{}{}}
\makeatother
\crefname{originallemma}{original lemma}{original lemmas}
\Crefname{originallemma}{Original lemma}{Original lemmas}
\makeatletter
\newlabel{morphisms-lemma-fppf-open}{{01UA}{source}{}{sourcepointer-4}{}}
\newlabel{morphisms-lemma-fppf-open@cref}{{[originallemma][0][]01UA}{[1][0][]source}{}{}{}}
\makeatother
\crefname{originallemma}{original lemma}{original lemmas}
\Crefname{originallemma}{Original lemma}{Original lemmas}
\makeatletter
\newlabel{schemes-lemma-specialize-points}{{01J7}{source}{}{sourcepointer-5}{}}
\newlabel{schemes-lemma-specialize-points@cref}{{[originallemma][0][]01J7}{[1][0][]source}{}{}{}}
\makeatother
\crefname{originaldefinition}{original definition}{original definitions}
\Crefname{originaldefinition}{Original definition}{Original definitions}
\makeatletter
\newlabel{topology-definition-Krull}{{0055}{source}{}{sourcepointer-6}{}}
\newlabel{topology-definition-Krull@cref}{{[originaldefinition][0][]0055}{[1][0][]source}{}{}{}}
\makeatother
\title{An etale morphism preserves the dimension of a scheme at a point}
\author{The Stacks Project Authors}
\date{stacks-project:descent.tex:local-snapshot-20261012; GNU FDL 1.2 or later}
\begin{document}
\maketitle
\noindent\textit{Editorial note.} All delivered mathematical text is the author's own, copied from the pinned Stacks Project chapter file descent.tex; no statement or proof step was rewritten, reordered or removed.\par\medskip
\noindent\textit{Editorial note.} Every reference of the proof points outside this chapter; no same-chapter context span is required and none was added.\par\medskip
\noindent\textit{Editorial note.} The counted spans carry 6 navigation stubs; no bibliography citation occurs in the spans.\par\medskip
\noindent\textit{Editorial note.} The preamble span points at the chapter's own preamble.tex shipped with the frozen source, used verbatim as staged.\par\medskip
\noindent\textit{Allowed standard prerequisites (editorial).}\begin{itemize}
\item Dimension at a point as the minimum of the Krull dimensions of the open neighbourhoods. Topology, Definition topology-definition-Krull, invoked at proof lines 5396-5399. Fixes the meaning of both sides of the asserted equality.
\item The fibres of an etale morphism over a field are discrete. Morphisms, Lemma morphisms-lemma-etale-over-field, invoked at proof lines 5412-5414. Guarantees that the images of two distinct members of a chain of specializations in U stay distinct in V.
\item Etale morphisms are open and flat, and a flat local ring map satisfies going down. Morphisms, Lemma morphisms-lemma-fppf-open at proof lines 5421-5422, Morphisms, Lemma morphisms-lemma-etale-flat at proof lines 5433-5434, and Algebra, Section algebra-section-going-up at proof lines 5437-5438. Openness keeps the image of a neighbourhood open, flatness makes the map of local rings flat and local, and going down lifts a chain of specializations along it.
\item Points of a scheme and specializations as primes of a local ring. Schemes, Lemma schemes-lemma-specialize-points, invoked at proof lines 5435-5436. Converts the chain of points in V into primes of the local ring, which is what going down applies to.
\end{itemize}\medskip
\ifdefined\hypertarget\hypertarget{sourcepointer-1}{}\fi
\noindent\textit{Source reference.} Original section 00HU: sources/stacks-chow/tags.txt, source lines 656--656. Pointer outside the counted span, invoked once for context; the load-bearing argument of the counted statement is the author's proof reproduced in full in the proof span.\par\medskip
\ifdefined\hypertarget\hypertarget{sourcepointer-2}{}\fi
\noindent\textit{Source reference.} Original lemma 02GS: sources/stacks-chow/tags.txt, source lines 3239--3239. Pointer outside the counted span, invoked once for context; the load-bearing argument of the counted statement is the author's proof reproduced in full in the proof span.\par\medskip
\ifdefined\hypertarget\hypertarget{sourcepointer-3}{}\fi
\noindent\textit{Source reference.} Original lemma 02GL: sources/stacks-chow/tags.txt, source lines 3232--3232. Pointer outside the counted span, invoked once for context; the load-bearing argument of the counted statement is the author's proof reproduced in full in the proof span.\par\medskip
\ifdefined\hypertarget\hypertarget{sourcepointer-4}{}\fi
\noindent\textit{Source reference.} Original lemma 01UA: sources/stacks-chow/tags.txt, source lines 2403--2403. Pointer outside the counted span, invoked once for context; the load-bearing argument of the counted statement is the author's proof reproduced in full in the proof span.\par\medskip
\ifdefined\hypertarget\hypertarget{sourcepointer-5}{}\fi
\noindent\textit{Source reference.} Original lemma 01J7: sources/stacks-chow/tags.txt, source lines 2004--2004. Pointer outside the counted span, invoked once for context; the load-bearing argument of the counted statement is the author's proof reproduced in full in the proof span.\par\medskip
\ifdefined\hypertarget\hypertarget{sourcepointer-6}{}\fi
\noindent\textit{Source reference.} Original definition 0055: sources/stacks-chow/tags.txt, source lines 197--197. Pointer outside the counted span, invoked once for context; the load-bearing argument of the counted statement is the author's proof reproduced in full in the proof span.\par\medskip
\begin{lemma}
\label{lemma-dimension-at-point-local}
Let $f : U \to V$ be an \'etale morphism of schemes.
Let $u \in U$ and $v = f(u)$. Then $\dim_u(U) = \dim_v(V)$.
\end{lemma}

\begin{proof}
In the statement $\dim_u(U)$ is the dimension of $U$ at $u$ as defined in
Topology, Definition \ref{topology-definition-Krull}
as the minimum of the Krull dimensions of open neighbourhoods of $u$ in $U$.
Similarly for $\dim_v(V)$.

\medskip\noindent
Let us show that $\dim_v(V) \geq \dim_u(U)$.
Let $V'$ be an open neighbourhood of $v$ in $V$.
Then there exists an open neighbourhood $U'$ of $u$ in $U$
contained in $f^{-1}(V')$ such that $\dim_u(U) = \dim(U')$. Suppose that
$Z_0 \subset Z_1 \subset \ldots \subset Z_n$ is a chain of irreducible
closed subschemes of $U'$. If $\xi_i \in Z_i$ is the generic point
then we have specializations
$\xi_n \leadsto \xi_{n - 1} \leadsto \ldots \leadsto \xi_0$.
This gives specializations
$f(\xi_n) \leadsto f(\xi_{n - 1}) \leadsto \ldots \leadsto f(\xi_0)$
in $V'$. Note that $f(\xi_j) \not = f(\xi_i)$ if $i \not = j$ as
the fibres of $f$ are discrete (see
Morphisms, Lemma \ref{morphisms-lemma-etale-over-field}).
Hence we see that $\dim(V') \geq n$. The inequality
$\dim_v(V) \geq \dim_u(U)$ follows formally.

\medskip\noindent
Let us show that $\dim_u(U) \geq \dim_v(V)$.
Let $U'$ be an open neighbourhood of $u$ in $U$.
Note that $V' = f(U')$ is an open neighbourhood of $v$ by
Morphisms, Lemma \ref{morphisms-lemma-fppf-open}.
Hence $\dim(V') \geq \dim_v(V)$. Pick a chain
$Z_0 \subset Z_1 \subset \ldots \subset Z_n$ of irreducible
closed subschemes of $V'$. Let $\xi_i \in Z_i$ be the generic point,
so we have specializations
$\xi_n \leadsto \xi_{n - 1} \leadsto \ldots \leadsto \xi_0$.
Since $\xi_0 \in f(U')$ we can find a point $\eta_0 \in U'$
with $f(\eta_0) = \xi_0$. Consider the map of local rings
$$
\mathcal{O}_{V', \xi_0} \longrightarrow \mathcal{O}_{U', \eta_0}
$$
which is a flat local ring map by
Morphisms, Lemma \ref{morphisms-lemma-etale-flat}.
Note that the points $\xi_i$ correspond to primes of the ring on the left by
Schemes, Lemma \ref{schemes-lemma-specialize-points}.
Hence by going down (see
Algebra, Section \ref{algebra-section-going-up})
for the displayed ring map we can find a sequence of specializations
$\eta_n \leadsto \eta_{n - 1} \leadsto \ldots \leadsto \eta_0$
in $U'$ mapping to the sequence
$\xi_n \leadsto \xi_{n - 1} \leadsto \ldots \leadsto \xi_0$
under $f$. This implies that $\dim_u(U) \geq \dim_v(V)$.
\end{proof}

\end{document}
